% !TeX root = main.tex
\begin{frame}{Zero-Knowledge Arguments of Knowledge (ZKAoK)}
	Originally proposed by Goldwasser, Micali and Rackoff \cite{GoldwasserMR85}.\pause
	
	Interactive communication between prover and verifier.\pause
	
	\begin{tikzpicture}
		\node[inner sep=0pt] (prover) at (0,0){\includegraphics[width=.2\textwidth]{prover.png}};
		\node[below = 0.3cm] at (prover.south) {Prover $(s, w)$};\pause
		\node[inner sep=0pt] (verifier) at (7,0) {\includegraphics[width=.2\textwidth]{verifier.png}};
		\node[below = 0.3cm] at (verifier.south) {Verifier $(s)$};\pause
		
		\draw[->,thick] (prover.north east) -- (verifier.north west) node[midway]{};
		\draw[->,thick] (verifier.west) -- (prover.east) node[midway]{};
		\draw[->,thick] (prover.south east) -- (verifier.south west) node[midway]{};
	\end{tikzpicture}
\end{frame} 
\begin{frame}{Zero-Knowledge Arguments of Knowledge (ZKAoK)}
	Properties:\pause
	\begin{itemize}
		\item {\small Completeness. Honest prover always succeeds.}\pause
		\item {\small Soundness. Cheating prover cannot convince verifier.}\pause
		\item {\small Zero-knowledge. Verifier gains nothing beyond validity of statement.}\pause
	\end{itemize}
	Zero-knowledge arguments of knowledge (ZKAoKs).\pause
	\begin{itemize}
		\item Argument: Computationally-bounded provers.\pause
		\item Argument of Knowledge: 
		
		$\exists$ extractor, with access to prover, to extract valid witness.\pause
	\end{itemize}

	Constructions in this thesis are ZKAoKs.\pause
	\begin{enumerate}[$\Rightarrow$]
		\item Privacy-preserving constructions.
	\end{enumerate}
\end{frame}